\section{The Resolver}
\label{sec:resolver}

The Zcash chain records candidate Name Notes. It does not identify which candidates were authorized by the Mint or which transition is current.

The Resolver derives this state from the canonical chain. It is permissionless. Two Resolvers that scan the same canonical chain under the same active manifest must derive the same registry.

\subsection{Scanning}

Standard Orchard wallets reject Name Notes because ZNS derives \(\psi\) and \(\mathsf{rcm}\) from the transition tuple rather than from the ordinary ZIP-212 \field{rseed} derivation. A Resolver uses a dedicated scanner and the published Name Note full viewing key.

For each candidate output, the Resolver performs three checks:

\begin{description}
  \item[Decryption.] The Resolver obtains the memo and note components by authenticated decryption, bypassing the standard ZIP-212 receiving check.
  \item[Name Note verification.] The Resolver parses the memo and reconstructs the Orchard commitment as specified in Section~\ref{sec:name-notes}.
  \item[Mint origin.] The Resolver accepts the candidate only if its creating Orchard action spends a note from the recognized Mint-controlled unspent set.
\end{description}

For every accepted Name Note, the Resolver updates the Mint-controlled unspent set and follows the \field{prev\_rcm} link for the affected name. When the best chain changes, it rolls back to the common ancestor and replays the replacement chain.

